In this chapter, we explore an alternative loss function, the ramp loss, inspired by the support vector machine (SVM).
The threshold model discussed in Section~\ref{ch:threshold-model} utilizes the so-called hard-margin loss (\autoref{fig:loss_functions}c), also known as \zo loss.
This means a correctly classified point induces a loss of $0$, while an incorrectly classified point induces a loss of $1$.
It may be beneficial to assign a lower loss to misclassified points that are ``close'' to the rule set.

The SVM is a mathematical programming-based binary classification method developed by Vapnik in 1998 and Cortes and Vapnik in 1995~\cite{vapnik_statistical_1998, cortes_support-vector_1995, brooks_support_2011}.
The classifier finds a hyperplane in the feature space that aims to separate positive datapoints from negative datapoints.
The formulation proposed by Vapnik and co-authors uses the hinge loss, a continuous measure for misclassification.
Specifically, the misclassification error is directly proportional to the distance between the point and the plane.
The behavior of the hinge loss is illustrated in \autoref{fig:loss_functions}a.
The hinge loss makes the SVM sensitive to outliers.
One method to increase robustness of the SVM, studied by several authors (\cite{brooks_support_2011, belotti_spatial_2024}),
is the use of the ramp loss, illustrated in \autoref{fig:loss_functions}b.

Besides being more robust than the hinge loss, the ramp loss has additional advantages.
Points that are ``close'' to a rule induce a lower loss.
This could reveal rules that would otherwise remain unchosen due to noise;
if many positive and negative points live near the border of an otherwise good rule,
the hard-margin loss can turn out large despite the rule being promising.
This is especially true in the case of a small selection of thresholds (i.e., a limited number quantile levels)
as opposed to full freedom in selecting thresholds.
In the former case, it may be that an available threshold happens to fall on the ``wrong side'' of a
relevant cloud of points.
A ramp loss could ease this pain by reducing the loss of points that are close to this threshold.
Additionally, the ramp loss could prove more robust than \zo loss,
as the change in loss as a result of a perturbation of the input is limited.
This is due to the reduced ``steepness'' of the ramp loss as opposed to the hard-margin loss.

This chapter covers the modification of the threshold model from the hard-margin loss to the ramp loss,
which is our contribution.
The resulting ramp loss model is a generalization of the threshold model.

Two variants are considered: ``additive'' and ``non-additive''.
The latter is an attempt to overcome issues with the former.
A mathematical description of the additive model is treated in
\autoref{sec:ramp-model-description} and \autoref{sec:ramp-pricing-problem}.
Examples of the model's behavior on simple, artificial examples are covered in \autoref{sec:ramp-examples}.
These examples are in two dimensions and help develop a visual intuition.
The subsequently developed non-additive variant is covered in \autoref{sec:ramp-nonadditive}.
Implementation is discussed in \autoref{sec:ramp-implementation}.
The experiments and their results are discussed
in \autoref{sec:ramp-methodology} and \autoref{sec:ramp-results}, respectively.

Readers who are mainly interested in the experiments and their results are advised to continue reading
\autoref{sec:ramp-examples}.
This is the section with examples, which helps to build an intuition for the model.
\autoref{sec:ramp-nonadditive} is mathematical and may be skipped by these readers,
but contains more examples worth reading about.

\begin{figure}
    \includesvg[width=0.5\textwidth]{gen/loss_functions_y_1_corrected.svg}
    \caption{Loss for a positive datapoint $i$ subject to rule $X_{\cdot j} \geq \alpha$ plotted
    against the feature value $X_{ij}$ with (a) the hinge loss, (b) the ramp loss and (c) the hard margin loss.
    The width of the margin is given by $\beta > 0$.
    An observation whose feature value $X_{ij}$ falls between $\alpha-\beta$ and $\alpha + \beta$ lies in the margin.}
    \label{fig:loss_functions}
\end{figure}


\section{Model description}\label{sec:ramp-model-description}
We will now reformulate \eqref{eq:BM} to incorporate a multidimensional generalization of the ramp loss.
To do this, first observe that the objective \eqref{eq:BM-obj} can be written as
\begin{equation*}
    \sum_{i \in \cP} \zeta_i + \sum_{i \in \cN} \sum_{k \in \K} \gamma_{ik} w_k,
\end{equation*} where
\[\gamma_{ik} = \begin{cases}
                    1, & \text{if $k \in \K_i$,} \\
                    0, & \text{otherwise}.
\end{cases}\]
In $\gamma_{ik}$, we recognize a hard-margin loss.
Specifically, for a negative point $i \in \cN$, the value $\gamma_{ik}$ is exactly the loss induced if $k$ is used.
Similarly, constraint \eqref{eq:BM_orig-misP} can be written as
\begin{equation*}
    \zeta_i + \sum_{k \in \K} \gamma_{ik} w_k \geq 1, \quad \forall i \in \cP.
\end{equation*}
If a positive point $i \in \cP$ is uncovered, this implies the sum will be zero and hence $\zeta_i = 1$, which corresponds to an induced loss of $1$.
The value $\gamma_{ik}$ represents a reduction of this loss if rule $k$ is used.
If $k$ covers the point, this value is $1$.
Using rule $k$ then allows $\zeta_i$ to be reduced to $0$.

By defining $\bm{\gamma}$ suitably, we can introduce a multidimensional variant of the ramp loss.
When doing so, the values of $\bm{\gamma}$ need no longer be binary, hence $\bm{\zeta}$ should no longer be binary.
Furthermore, the ramp loss of a point is anywhere between $0$ and $2$.
The ramp loss model is therefore given by
\begin{subequations}
    \label{eq:RM}
    \begin{align}
        \mathbf{\min} \quad & \sum_{i \in \cP} \zeta_i + \sum_{i \in \cN} \sum_{k \in \K} \gamma_{ik} w_k \label{eq:RM-obj} \\
        \mathbf{\text{s.t.}} \quad & \zeta_i + \sum_{k \in \K} \gamma_{ik} w_k \geq 2, & \forall i \in \cP \label{eq:RM-misP} \\
        & \sum_{k \in \K} c_k w_k \leq C, & \label{eq:RM-complex} \\
        & w \in \{0, 1\}^{\K}, \zeta \in [0,2]^{\cP}, \label{eq:RM-vars}
    \end{align}
\end{subequations}
To define $\bm{\gamma}$, we first define a decision function $\bm{d}: \P(\L) \to \R^\I$ such that $d_i(k)$ is
non-negative if rule $k$ covers point $i$ and negative otherwise.
Furthermore, the value $d_i(k)$ quantifies how far $i$ is from the boundary of $k$.
Define
\begin{equation}\label{eq:decision_function}
    d_i(k) = \umin{(j,\alpha) \in k} X_{ij} - \alpha, \tquad i \in \I, \; k \in \K\setminus\{\emptyset\}.
\end{equation}
If $k$ covers $i$, then any element in the minimum is non-negative, meaning $d_i(k)$ is the distance to the closest threshold.
If rule $k$ does not cover $i$, then there exists some clause $(j,\alpha) \in k$ that cuts off $i$.
Choose such $(j,\alpha)$.
It follows that $X_{ij} - \alpha < 0$ and so $d_i(k) < 0$.
By extension, we have in this case that
\[
    d_i(k) = -\umax{(j,\alpha) \in k \fw X_{ij} < \alpha} \alpha - X_{ij}, \tquad i \in \I.
\]
This means that $d_i(k)$ is the (negated) distance to the furthest unsatisfied threshold.
Therefore, the value $-d_i(k)$ can be interpreted as a distance to the covered region.
To account for the empty rule, we set
\begin{equation}
    d_i(\emptyset) = \infty, \tquad i \in \I. \label{eq:decision_function_empty_rule}
\end{equation}
Denote by $\beta$ the margin width, as illustrated in \autoref{fig:loss_functions}b.
While the margin width is a variable in the Support Vector Machine~\cite{belotti_spatial_2024},
it is a constant in our model.
We define
\begin{equation}\label{eq:gamma_ik}
    \gamma_{ik} = \begin{cases}
                      0, & \text{if $d_i(k) \leq -\beta$}, \\
                      1 + \frac{1}{\beta}d_i(k), & \text{if $d_{i}(k) \in (-\beta, \beta]$}, \\
                      2, & \text{if $d_i(k) > \beta$},
    \end{cases} \tquad i \in \I, \; k \in \K.
\end{equation}
This definition realizes a generalized ramp loss in the master problem \eqref{eq:RM}.

We refer to this model as the ``additive'' ramp loss model.
This name is due to the additive nature of loss reduction for positive points.
Recall \eqref{eq:RM-misP}.
Several rules, when used together, can all contribute to reducing a $\zeta$-variable.
For example, if point $i$ lies in several margins, the corresponding loss $\zeta_i$ can still be $0$,
which is equivalent to being far inside the covered region of the ruleset.


As $\beta \to 0,$ the ramp loss converges pointwise to the hard-margin loss (multiplied by $2$)
for all values of $X_{ij}$ except the corresponding $\alpha$.
Assuming no feature values lies on such threshold,
it follows that the ramp loss model is equivalent to the threshold model when $\beta$ is arbitrarily small.
To see this, observe that for $\beta$ small enough, $\bm{\gamma}$ and $\bm{\zeta}$ are those used in
the threshold model, multiplied by $2$.

In a realistic setting, it is to be expected that few, or none, of the datapoints lie exactly on a threshold.
There are two reasons for this.
First, continuous measurements typically include a measurement error,
making it less likely for several datapoints to have an equal feature value.
Second, the linear interpolation used in the quantile function is likely to yield a non-observed value.
Therefore, in a practical setting with continuous features,
the ramp loss model with $\beta$ arbitrarily small is practically equivalent to the threshold model.
For convenience, we define for the special case $\beta = 0$,
\begin{equation}\label{eq:ramp_gamma_mw_0}
    \begin{split}
        \gamma_{ik} &= 2 \cdot \Ind(d_i(k) \geq 0) \\
        &=
        \begin{cases}
            0, & \text{if $\exists (j,\alpha)\in k \fw X_{ij} < \alpha$},\\
            2, & \text{if $X_{ij} \geq \alpha \quad \forall (j,\alpha) \in k$},\\
        \end{cases} \tquad i \in \I, \; k \in \K.
    \end{split}
\end{equation}
For this choice of $\bm{\gamma}$, the model is exactly equivalent to the threshold model.
We then use the implementation of the ramp loss model with $\beta = 0$ to solve the threshold model.
The goal is now to express the pricing problem in terms of $\bm{\gamma}$ and study how this problem may be solved.

\section{Pricing problem}\label{sec:ramp-pricing-problem}
The reformulations \eqref{eq:RM-obj} and \eqref{eq:RM-misP} affect the pricing problem.
We will now derive the pricing problem, study its complexity and construct a mixed integer program (MIP) that can be used to solve it.
Given non-negative dual values $(\bm{\mu}, \lambda) \in \R_+^\P \times \R_+$, the reduced cost of a rule $k \in \K$ is given by
\begin{equation} \label{eq:RM-redcost}
    \rho_k = \sum_{i \in \cN} \gamma_{ik} - \sum_{i \in \cP} \mu_{i}\gamma_{ik} + \lambda c_k.
\end{equation}
The derivation is discussed in \autoref{sec:reduced_costs_additive_ramp_loss}.
Denote by $[\cdot]_a^b$ the clipping operator that clamps the value of $\cdot$ to $[a,b].$
For $\beta > 0$, we have that
\begin{equation}\label{eq:gamma_ik_min}
\begin{split}
        \gamma_{ik} &= \left[1 + \frac{1}{\beta} d_i(k) \right]_0^2 \\
        &= \left[1 + \frac{1}{\beta}\umin{(j,\alpha) \in k} X_{ij} - \alpha\right]_0^2 \\
        &= \umin{(j,\alpha) \in k} \left[1 + \frac{1}{\beta}(X_{ij} - \alpha)\right]_{0}^{2}, \tquad i \in \I,\; k \in \K\setminus\{\emptyset\}.
\end{split}
\end{equation}
Define
\begin{equation}\label{eq:gamma_{il}}
    \gamma_{il} = \begin{cases}
                      \left[1 + \frac{1}{\beta}(X_{ij} - \alpha)\right]_{0}^{2}, & \text{if $\beta > 0$,} \\
                      2\cdot \Ind(X_{ij} \geq \alpha), & \text{if $\beta = 0$,}
    \end{cases} \tquad i \in \I, \; l = (j,\alpha) \in \L.
\end{equation}
Then
\begin{equation}\label{eq:f_min}
    \gamma_{ik} = f_i(k) \coloneq \begin{cases}
                               2, & k = \emptyset, \\
                               \umin{l \in k} \gamma_{il}, & k\neq \emptyset, \\
    \end{cases} \tquad i \in \I, k \in \K.
\end{equation}
Minimizing \eqref{eq:RM-redcost} is equivalent to solving
\begin{equation}\label{eq:RP_set_formulation}
    \umin{S \subseteq \L} \quad \sum\limits_{i \in \N} f_i(S)  - \sum\limits_{i \in \P}  \mu_i f_i(S) + \lambda c(S),
\end{equation}
where $c: \K \to \Z_+$ is a function measuring rule complexity.

%Second, we may restrict the use of a clause if a more restrictive clause of the same feature is already chosen.
%To describe this, order for each $j \in \J$ the clauses involving feature $j$ by the respective threshold $\alpha$ in non-decreasing order.
%denote by $(j,(\eta))$ the clause involving feature $j$ with the $\eta$-th largest threshold.
%Then the constraint is given by
%\[z_{(j,(\eta))} \leq (1-)\]
%Note that these two constraints have an identical effect for the binary problem.
%However, for the LP relaxation, the constraints are not equivalent.

\subsection{Complexity}\label{subsec:ramp-pricing-problem-complexity}
In this section we discuss the complexity of the pricing problem \eqref{eq:RP_set_formulation} for various definitions of rule complexity.
It turns out to be NP-hard in general.
In \autoref{ch:boolean-model}, the rule complexity of a rule $S \subseteq \L$ was defined as \[c(S)=1+|S|.\]
\begin{theorem}
    \cite{lawless_interpretable_2023}
    \label{thm:np-hard-lawless}
    The problem \eqref{eq:RP_set_formulation} is NP-hard for $c(S) = 1 + |S|.$
\end{theorem}

\begin{proof}
    The proof is analogous to the one given in~\cite{lawless_interpretable_2023} and is by reduction from the minimum vertex cover (MVC) problem, which is known to be NP-complete.
    Let $G=(V,E)$ be an undirected graph with at least one edge, otherwise the problem has a trivial solution.
    We construct an instance of \eqref{eq:RP_set_formulation} that finds a minimum cardinality vertex cover of $G$.
    Take $\P = \emptyset.$
    For each vertex $v \in V,$ create a clause $l_v = (v, 1) \in \L$.
    For each edge $e = \{u,v\} \in E,$ create a point $i_e \in \N$ such that
    \[X_{i_e l_v} = \begin{cases}
                        0, & \text{if $v \in e$}, \\
                        2, & \text{otherwise}.
    \end{cases}\]
    Choose $\beta = 1.$ It follows that
    \[
        f_{i_{e}}(k) \in \{0, 2\}, \tquad e \in E, \; k \in \K.
    \]
    Precisely, we have that
    \[f_{i_e}(S) = \begin{cases}
                         0, & \text{if $\exists l_v \in S \fw v \in e$,} \\
                         2, & \text{otherwise},
    \end{cases} \tquad i_e \in \N, \; S \subseteq \L.
    \]
    Additionally, let $\lambda = 1/(1 + |\L|).$
    This leads to the problem
    \[\umin{S \subseteq \L} \quad \sum\limits_{i \in \N} f_i(S) + \frac{1 + |S|}{1 + |V|}.\]
    If $S$ is not a vertex cover, then there is some $e \in E$ such that $f_{i_e}(S) = 2.$ So
    \[
        \sum\limits_{i \in \N} f_i(S) \geq 2.
    \]
    Consequently, the objective value is at least $2$.
    If $S$ is a vertex cover, then the sum evaluates to $0$, since for all $e \in E$ we have $f_{i_e}(S) = 0.$
    In this case, the objective is at most $1$, because
    \[\frac{1 + |S|}{1 + |V|} \leq 1.\]
    Since this objective value is strictly better,
    an optimal solution $S$ is a vertex cover and minimizes the above fraction.
    This is equivalent to minimizing $|S|$.
    So the problem finds a minimum cardinality vertex cover.
    Therefore, this instance of the pricing problem, constructible in polynomial time, solves the MVC problem. \qedhere
\end{proof}

Analogously, the problem \eqref{eq:RP_set_formulation} is NP-hard for $c(S) = |S|.$
This proof is not applicable for the case $c(S) = 1,$ which was chosen in \autoref{ch:threshold-model}.

\begin{theorem}
    \label{thm:np-hard}
    The problem \eqref{eq:RP_set_formulation} is NP-hard for $c(S) = 1$.
\end{theorem}

\begin{proof}
    The proof is again by reduction from the minimum vertex cover (MVC) problem and uses a similar construction.
    Let $G=(V,E)$ be an undirected graph with at least one edge.
    Recall the instance described in the proof of \autoref{thm:np-hard-lawless}.
    Augment the instance by creating for each $v \in V$ a point $i_v \in \P$ with
    \[X_{i_v l_{\overline{v}}} = \begin{cases}
                                    0, & \text{if $v = \overline{v}$}, \\
                                    2, & \text{otherwise},
    \end{cases} \tquad v,\overline{v} \in V.\]
    It follows that
    \[
        f_{i_v}(S) = \begin{cases}
                           0, & \text{if $l_v \in S$,} \\
                           2, & \text{otherwise},
        \end{cases}
        \tquad v \in V, \; S \subseteq \L.
    \]
    Let $\lambda = 0$ and $\mu_i = 1/(1 + |V|), \; i \in \P.$
    The problem is then given by
    \[\umin{S \subseteq \L} \quad \sum\limits_{i \in \N} f_i(S) - \frac{1}{1+|V|} \sum\limits_{i \in \P} f_i(S).\]
    Observe that the solution containing every $l_v, \; v \in V$ has objective value $0$.
    Let $S$ be an optimal solution.
    Then the corresponding objective value is at most $0$, which implies
    \[
        \sum\limits_{i \in \N} f_i(S) \leq \frac{1}{1+|V|} \sum\limits_{i \in \P} f_i(S)
    \]
    Note that
    \[
        \frac{1}{1+|V|} \sum\limits_{i \in \P} f_i(S) \leq \frac{1}{1+|V|} \sum\limits_{i \in \P} 2 = \frac{2|V|}{1+|V|} < 2,
    \]
    Therefore
    \[
        \sum\limits_{i \in \N} f_i(S) < 2.
    \]
    Since $f_i(S) \in \{0,2\}$ for all $i$,
    it follows that
    \[\sum\limits_{i \in \N} f_i(S) = 0.\]
    This again implies that $S$ corresponds to a vertex cover.
    It also follows that solution $S$ minimizes
    \[
        \begin{split}
            -\frac{1}{1+|V|} \sum\limits_{i \in \P} f_i(S) &= -\frac{1}{1+|V|} \sum\limits_{i_v \in \P} 2 \cdot \Ind(l_v \notin S)
            \\ &= -\frac{2\left(|V| - |S|\right)}{1 + |V|} \\ &= -\frac{2|V|}{1 + |V|} + \frac{2|S|}{1+|V|}.
        \end{split}
        \]
    which is equivalent to minimizing $|S|$.
    Therefore, this instance of the pricing problem, constructible in polynomial time, solves the MVC problem. \qedhere
\end{proof}


\subsection{Mixed Integer Programming formulations}\label{subsec:ramp-pricing-problem-mip}
In this section, we explore a MIP formulation of the pricing problem \eqref{eq:RP_set_formulation} and discuss challenges in solving it.
Recall the reduced cost \eqref{eq:RM-redcost}.
Let $\bm{z}\in \{0,1\}^{\L}$ again be integer variables representing which clauses are used in the rule.
Consider the MIP
\begin{subequations}
    \label{eq:RP}
    \begin{align}
        \min \quad & \sum_{i \in \cN} \gamma_i - \sum_{i \in \cP} \mu_i \gamma_i +
        \lambda(\bm{z}) \label{eq:RP-obj} \\
        \st \quad & \gamma_i \leq 2 - (2 - \gamma_{il})z_l, & l \in \L, &\; i \in \cP, \label{eq:RP-pos} \\
        & \gamma_i \geq \gamma_{il} - \sum_{l' \in \L \fw \gamma_{il'} < \gamma_{il}} (\gamma_{il} - \gamma_{il'}) z_{l'} , & l \in \cL, &\; i \in \cN, \label{eq:RP-neg} \\
        & \gamma_i \geq 2 - \sum_{l \in \L} 2z_l, & l \in \L, \label{eq:RP-fix} \\
        & \sum_{l \in \L} z_l \leq D, \label{eq:RP-complex}\\
        & \bm{z}\in \{0,1\}^{\L},~ \bm{\gamma}\in [0,2]^\I.\label{eq:RP-vars}\qquad
    \end{align}
\end{subequations}
We also impose the one-clause-per-feature constraint as given in \eqref{eq:TP_one-per-feature}.

In this formulation, constraints \eqref{eq:RP-pos} ensure that $\gamma_i$ for positive $i$ reflects $\gamma_{ik},$
where $k$ is the chosen rule.
Strictly speaking, this is only true for $i \in \P$ for which $\mu_i > 0$.
For $i \in \P$ such that $\mu_i = 0$, the variable $\gamma_i$ is effectively redundant,
so these variables need not be considered.
For negative $i$, constraints \eqref{eq:RP-neg} and \eqref{eq:RP-fix} ensure this.
We will now prove the claim for both positive and negative points.

Denote by $\gopt, \zopt$ an optimal solution for this IP\@.
Recall that $\zopt$ is binary.
Define
\[\sopt = \{l \in \L \fw \zopt[l] = 1\},\]
and recall the definition of $f$ in \eqref{eq:f_min}.

\begin{claim}\label{claim:TP_ramp-consP}
    For all $i \in \P$ such that $\mu_i > 0,$ it holds that
    \[\gopt[i] = f_i(\sopt).\]
\end{claim}

\begin{proof}
    Let $i \in \P$ such that $\mu_i > 0$ and let $l \in \L$.
    If $\zopt[l] = 0$, the corresponding constraint \eqref{eq:RP-pos} becomes redundant.
    If $\zopt[l] = 1,$ the constraint becomes $\gamma_i \leq \gamma_{il}.$
    Therefore, if $\sopt \neq \emptyset$ it holds that
    \[\gamma_i \leq \umin{l \in \sopt} \gamma_{il}.\]
    If $\sopt = \emptyset$, we have that $\gamma_i \leq 2$ by the variable bound.
    In any case,
    \[
        \gamma_i \leq f_i(\sopt).
    \]
    There are no other constraints involving $\gamma_i$.
    Since $\gamma_i$ has a negative coefficient in the objective, this inequality will be an equality in an optimal solution.
    \qedhere
\end{proof}

\begin{claim}\label{claim:TP_ramp-consN}
    For all $i \in \N$ it holds that
    \[\gopt[i] = f_i(\sopt).\]
\end{claim}


\begin{proof}
    Let $i \in \N$.
    If $\sopt = \emptyset,$ then $\gopt[i] \geq 2 = f_i(\emptyset)$ by constraint \eqref{eq:RP-fix}.
    So if $\sopt = \emptyset$, the statement holds.

    Assume $\sopt \neq \emptyset.$
    Observe that constraint $\eqref{eq:RP-fix}$ becomes redundant, as there is at least one $l$ such that $z_l = 1$.
    Let $l^* \in \sopt$ such that \[\gamma_{il^*} = \umin{l \in \sopt} \gamma_{il}.\]
    For this $i$ and $l$, the constraint \eqref{eq:RP-neg} is
    \[\gamma_i \geq \gamma_{il}.\]
    Now let $l \in \L$.
    We consider the right hand side of \eqref{eq:RP-neg}.
    We have that
    \[\begin{split}
          \gamma_{il} - \sum_{l' \in \L\,:\,\gamma_{il'} < \gamma_{il}} (\gamma_{il} - \gamma_{il'}) z_l
          &= \gamma_{il} - \sum_{l' \in \sopt \,:\,\gamma_{il'} < \gamma_{il}} (\gamma_{il} - \gamma_{il'}).
    \end{split}\]
    If the sum is empty, then it must hold that $\gamma_{il} = \gamma_{il^*}$ and the right hand side equals $\gamma_{il^*}$.
    If the sum is non-empty, we have
    \[
        \begin{split}
            \gamma_{il} - \sum_{l' \in \sopt \,:\,\gamma_{il'} < \gamma_{il}} (\gamma_{il} - \gamma_{il'}) &\leq \gamma_{il} - \umax{l' \in \sopt \,:\,\gamma_{il'} < \gamma_{il}} (\gamma_{il} - \gamma_{il'}) \\
            &= \umin{l' \in \sopt \,:\,\gamma_{il'} < \gamma_{il}} \gamma_{il'} \\
            &= \gamma_{il^*}.
        \end{split}
    \]
    Hence, $\gamma_i \geq \gamma_{il^*} = \umin{l \in \sopt} \gamma_{il}$ is the tightest constraint on $\gamma_i$, induced by $l^*$.

    Combining the case $\sopt = \emptyset$ and $\sopt \neq \emptyset,$ we get that
    \[
        y_i \geq f_i(\sopt),
    \]
    is the tightest lower-bound on $\gamma_i$.
    Since $\gamma_i$ has a positive coefficient in the objective, this inequality will be an equality for an optimal solution.
    From this, the result follows.
\end{proof}

\begin{note}
    In the process of constructing and implementing this MIP, constraint \eqref{eq:RP-fix} was overlooked.
    Hence, it was not implemented when the experiments in this report were conducted.
    Luckily, this mistake does not affect the experiments for various reasons.
    First, the solving the MIP could not be done in reasonable time for the PD-L1 experiments,
    hence only the unaffected heuristic pricer was used.
    If the exact pricer implementation without constraint \eqref{eq:RP-fix} is used anyway, the effect is still limited.
    We note that for any non-zero solution $\bm{z} \in \{0,1\}^\L,$ the variables $\bm{\gamma}$ are correctly modeled.
    Hence, the reduced cost can only be incorrect when $\bm{z} = 0$.
    For $\bm{z} = 0$, it can happen that some $\gamma_i$ are too low for negative $i$.
    If so, the reduced cost is underestimated.
    This is only problematic if, due to this error, the empty rule becomes the rule with the lowest reduced cost.
    In this case, the rule would not be added to the master problem in our implementation,
    and pricing would continue (with new shadow prices),
    which may cause a rule with true negative reduced cost to be found.
    However, the problem of underestimating $\gamma_i$ for negative $i$ only occurs if there is no
    clause $l$ such that $\gamma_{il} = 2$.
    This is the case when there is no clause that covers (further than the margin) the negative point.
    In practice, we do not expect this to happen often,
    as we include upper-bounds with low thresholds and lower-bounds with high thresholds.
\end{note}

The LP relaxation of \eqref{eq:RP} is not tight.
The LP relaxation allows for \emph{spreading}, by which several entries of $\bm{z}$ admit a fractional value.
\begin{example}[Spreading]\label{ex:branch1}
    \begin{figure}
        \includesvg[width=0.5\textwidth]{gen/branch_ex1.svg}
        \caption{
            An example problem with a two-dimensional feature space, composed of one positive datapoint $\bm{X}_2 = (2\tfrac{1}{2}, 2\tfrac{1}{2})$, two negative datapoints $\bm{X}_1 = (1\tfrac{1}{2}, 1\tfrac{1}{2}), \bm{X}_3 = (3\tfrac{1}{4}, 3\tfrac{1}{4})$ and clauses $a$, $b$, $f$ representing $X_{\cdot 1} \geq 3$,$X_{\cdot 2} \geq 3$ $X_{\cdot 1} \geq \tfrac{1}{4}$, respectively.
            The margin width $\beta=1.$
            The boundary of the margin is indicated with dotted lines.}
        \label{fig:branch1}
    \end{figure}
    Let
    %
    \[
        \mathcal{N} = \{1, 3\}, \quad \mathcal{P} = \{2\}, \quad
        \mathcal{L} = \{a, b, f\}.
    \]
    Consider the two-dimensional feature space and lower-bounds in \autoref{fig:branch1}.
    The matrix $(\gamma_{il})$ is given by
    %
    \[
        \begin{array}{c|ccc}
            & a    & b    & f \\
            \hline
            1 & 0    & 0    & 1 \\
            2 & 1/2  & 1/2  & 1 \\
            3 & 1/4  & 1/4  & 1 \\
        \end{array}
    \]
    Choose $D = 3.$ In the initial pricing round, $\mu_i = 1$ and $\lambda = 0$.
    The LP relaxation of~\eqref{eq:RP} attains an optimal objective value of $-1/2,$
    by
    %
    \[
        \bm{z} = \bigl(\tfrac{1}{2},\, \tfrac{1}{2},\, 0\bigr), \qquad
        \gamma_1 = 0, \quad \gamma_2 = \tfrac{3}{4}, \quad \gamma_3 = \tfrac{1}{4}.
    \]
    Given $\bm{z}$, the values of $\bm{\gamma}$ can be computed directly.
    The optimal objective value of the IP is $-1/4, $ attained by several solutions, examples being
    \[\bm{z}=(1,0,0), \quad \bm{z}=(0,1,0), \quad \bm{z}=(1,1,0).\]
    For any of these solutions, it holds that $\gamma_2 \leq 1/2$ because at least one of $a$ or $b$ is used.
    Meanwhile, we have $\gamma_2 = \frac{3}{4}$ in an optimal LP solution, like the one shown above.
    We refer to this as spreading, meaning the use of fractional values in $\bm{z}$ to admit larger $\gamma_i$ for $i \in \P,$ leading to a lower objective value.

    In this example, the issue is resolved when branching on $z_a$ or $z_b$.
    Under any of the branching decisions
    \[z_a = 0,\; z_a = 1,\; z_b = 0, \text{ or } z_b = 1,\]
    the resulting LP solution is integer and $\gamma_2 = 1/2.$
\end{example}

%\begin{example}[Spreading and branching]\label{ex:branch2}
%    \begin{figure}
%        \includesvg[width=0.5\textwidth]{gen/branch_ex2.svg}
%        \caption{
%            An example problem with a two-dimensional feature space, composed of one positive datapoint $\bm{X}_2 = (2\tfrac{1}{2}, 2\tfrac{1}{2})$, two negative datapoints $\bm{X}_1 = (1\tfrac{1}{2}, 1\tfrac{1}{2}), \bm{X}_3 = (3\tfrac{1}{4}, 3\tfrac{1}{4})$ and clauses $a, b, c, d, f$, representing $X_{\cdot 1} \geq 3$, $X_{\cdot 2} \geq 3$, $X_{\cdot 2} \geq 2\tfrac{3}{4}$, $X_{\cdot 2} \geq 2\tfrac{3}{4}$ and $X_{\cdot 1} \geq \tfrac{1}{4}$, respectively.
%            The margin width $\beta=1.$
%            The boundary of the margin is not shown.}
%        \label{fig:branch2}
%    \end{figure}
%
%    We consider the problem discussed in \autoref{ex:branch1} with two additional clauses $c$ and $d$, see \autoref{fig:branch2}.
%    The matrix $(\gamma_{il})$ is given by
%    %
%    \[
%        \begin{array}{c|ccccc}
%              & a    & b    & c    & d    & f \\
%            \hline
%            1 & 0    & 0    & 0    & 0    & 1 \\
%            2 & 1/2  & 1/2  & 3/4  & 3/4  & 1 \\
%            3 & 1/4  & 1/4  & 1/2  & 1/2  & 1 \\
%        \end{array}
%    \]
%    The LP relaxation of~\eqref{eq:RP} attains an optimal objective value of $-1/2,$
%    achieved for example by
%    %
%    \[
%        \bm{z} = \bigl(\tfrac{1}{4},\, \tfrac{1}{4},\, \tfrac{1}{4},\,
%        \tfrac{1}{4},\, 0\bigr), \qquad
%        \gamma_1 = 0, \quad \gamma_2 = \tfrac{7}{8}, \quad \gamma_3 = \tfrac{3}{8}.
%    \]
%    %
%    In any solution to the IP that uses clause $a$ or $b$, it holds that $\gamma_2 \leq 1/2.$
%    In any solution that uses clause $c$ or $d$, it holds that $\gamma_2 \leq 3/4.$
%    Spreading is not necessarily prevented by branching on a variable with fractional value.
%    Suppose we branch on $z_d.$ If $z_d = 0$ is fixed, there exists optimal LP solution
%    \[
%        \bm{z} = \bigl(\tfrac{1}{4},\, \tfrac{1}{4},\, \tfrac{1}{4},\,
%        \tfrac{1}{2},\, 0\bigr), \qquad
%        \gamma_1 = 0, \quad \gamma_2 = \tfrac{7}{8}, \quad \gamma_3 = \tfrac{3}{8},
%    \]
%    with objective value $-1/2.$
%    If $z_d = 1$ is fixed, an optimal solution is given by
%    \[
%        \bm{z} = \bigl(\tfrac{1}{2},\, \tfrac{1}{2},\, \tfrac{1}{4},\,0,\, 1\bigr), \qquad
%        \gamma_1 = 0, \quad \gamma_2 = \tfrac{3}{4}, \quad \gamma_3 = \tfrac{1}{4},
%    \]
%    again with objective value $-1/2.$
%    If constraint \eqref{eq:TP_one-per-feature} is also added, an optimal solution when $z_d = 1$ is fixed is given by
%    \[
%        \bm{z} = \bigl(\tfrac{1}{2},\, 0,\, \tfrac{1}{2},\,1,\, 1\bigr), \qquad
%        \gamma_1 = 0, \quad \gamma_2 = \tfrac{3}{4}, \quad \gamma_3 = \tfrac{3}{8},
%    \]
%    and has objective value $-3/8.$
%    In all cases, an optimal LP solution after branching is still fractional.
%    An analogous result is observed when branching on $z_c$.
%    Branching on $z_a$ or $z_b$ is better.
%    If $z_a = 1$ is fixed, for example, an optimal LP solution has objective value $-1/4$, the optimal IP objective value.
%    The solution found by SCIP is integer, namely the solution that uses clauses $a$ and $b$.
%    If $z_a = 0$ is fixed, the optimal LP objective value is $-9/20$, attained for example by
%    \[
%        \bm{z} = \bigl(0,\, \tfrac{1}{5},\, \tfrac{2}{5},\,
%        \tfrac{2}{5},\, 0\bigr), \qquad
%        \gamma_1 = 0, \quad \gamma_2 = \tfrac{9}{10}, \quad \gamma_3 = \tfrac{9}{20}.
%    \]
%    \todo{How would SCIP choose a variable to branch on? Would it systematically avoid $z_d$ and $z_c$?
%    How would it work in more complicated examples? Is it then possible to choose these 'useless' branching variables?}
%
%    An alternative approach is branching on a set of variables.
%    For instance, choose $S=\{a, b\}$ and consider the two branches
%    \[\sum_{l \in S} z_l \geq 1, \text{ and } \sum_{l \in S} z_l = 0.\]
%    This branching strategy by itself is not more effective than variable branching.
%    However, each branch induces problem-specific variable bounds not captured by the LP\@.
%    In the former branch, at least one of $a$ or $b$ is chosen, and therefore \[\gamma_2 \leq 1/2.\]
%    Adding this constraint, which is equivalent to an updated upper bound on $\gamma_2$, this branch yields an integer solution.
%    In the latter branch, it follows that \[\gamma_3 \geq 1/2.\]
%    Even with this induced constraint, the branch does not yield an integer solution.
%    Instead, spreading is observed over $c$ and $d$.
%    Note that the set branching is a generalization of variable branching, as branching on some $S_1$ with one element is equivalent to branching on a variable.
%\end{example}

%\subsection{Custom branching rule}\label{subsec:ramp-branching-rule}
%
%
%We now introduce the problem-specific set branching approach.
%Each node $\nu$ corresponds to a set $\bs^\nu$, which contains disjoint subsets of $\L$.
%In each node $\nu,$ we restrict the problem to
%\begin{equation}\label{eq:custom_branching_set_constraints}
%    \forall_{S \in \bs^\nu} \sum_{l \in S} z_l \geq 1, \quad \text{and} \quad \forall_{l \notin \bigcup_{S \in \bs^\nu} S} \; z_l = 0.
%\end{equation}
%Furthermore, we may update bounds on $\gamma_i$ for all $i \in \pun.$ Precisely,
%\begin{align}
%    \gamma_i & \leq \underset{S \in \bs^\nu}{\min} \underset{l \in S}{\max} \gamma_{il} ,  & i \in \P, \label{eq:custom_branching_bound_P} \\
%    \gamma_i & \geq \underset{l \in \bigcup_{S \in \bs^\nu}}{\min} \gamma_{il} ,  & i \in \N. \label{eq:custom_branching_bound_N} \\
%\end{align}
%Constraint \eqref{eq:custom_branching_bound_P} upper bounds $\gamma_i$,
%using the fact that at least one clause per $S \in \bs^\nu$ will be chosen.
%While this constraint is valid, a possibly more powerful approach is to change constraints \eqref{eq:TP_ramp-pos}.
%To do this, define $S^\nu_i$, $i \in \P$, as a set such that
%\[\underset{l \in S^\nu}{\max} \gamma_{il} = \underset{S \in \bs^\nu}{\min} \underset{l \in S}{\max} \gamma_{il}.\]
%Essentially, the set $S^\nu_i$ is the set with the tightest $\gamma_i$-bound over all $S \in \bs^\nu$.
%Then define
%\[\hat{\gamma}^\nu_i  =  \underset{l \in S^\nu_i}{\max} \gamma_{il}. \]
%Now replace constraints \eqref{eq:TP_ramp-pos} by
%\[ \gamma_i \leq \hat{\gamma}^\nu_i- (\hat{\gamma}^\nu_i - \gamma_{il})z_l, \quad i \in \cP, \; l \in \L \text{ such that } \gamma_{il} < \hat{\gamma}^\nu_i. \label{eq:TP_ramp-pos_branch} \\\]
%Constraint \eqref{eq:custom_branching_bound_N} lower-bounds $\gamma_i$,
%using the fact that only clauses inside some set $S \in \bs^\nu$ can be used.
%However, this constraint follows directly from constraints \eqref{eq:TP_ramp-neg} when $\forall_{l \notin \bigcup_{S \in \bs^\nu}} \; z_l = 0$.
%The branching procedure is as follows:
%\begin{enumerate}
%    \item At node $\nu$, choose a set $S \in \bs^\nu$ such that $|S| \geq 2$,
%    \item Partition $S$ into subsets $S_1, S_2$,
%    \item Create the following branches:
%    \begin{enumerate}
%        \item A branch $\nu_1$ with $\bs^{\nu_1} = \bs^\nu \cup \{S_1\} \setminus{\{S\}},$
%        \item A branch $\nu_2$ with $\bs^{\nu_2} = \bs^\nu \cup \{S_2\} \setminus{\{S\}},$
%        \item If $|\bs^\nu| < D,$ a branch $\nu_3$ with $S^{\nu_3} = \bs^\nu \cup \{S_1, S_2\} \setminus{\{S\}}.$
%    \end{enumerate}
%\end{enumerate}
%If no more branching is possible at a node $\nu$, meaning step 1 fails, then $\bs^\nu$ contains only sets of size 1.
%Hence, it encodes a binary solution.
%The root node $\nu_\text{root}$ has $\bs^{\nu_\text{root}} = \{\L\}$.
%
%The idea is that clever choices in steps 1 and 2 of the branching procedure yield a binary solution high up in the branching tree,
%thanks to the continuously updated bounds on $\bm{\gamma}$.
%This raises the question: how should $S$ and $S_1, S_2$ be chosen at each node?
%
%Denote the optimal LP solution at node $\nu$ by $\overline{\bm{z}}^\nu, \overline{\bm{\gamma}}^\nu.$
%Define
%\[T^\nu = \{l \in \L \fw \overline{z}^\nu_l > 0 \}. \]
%Then define
%\[T_i^\nu = \{l \in T^\nu \fw  \overline{\gamma}^\nu_i > \gamma_{il}\}, \quad i \in \P. \]
%Observe that for any $l \in T^\nu,$ we have by constraint \eqref{eq:TP_ramp-pos} that
%\[\overline{z}^\nu_l = 1 \implies \overline{\gamma}^\nu_i \leq \gamma_{il} \iff l \notin T_i^\nu. \]
%Conversely,
%\[l \in T_i^\nu \implies \overline{z}^\nu_l \in (0,1). \]
%Hence, if there exists $i \in P$ such that $T_i^\nu \neq \emptyset,$ variables in this set are good candidates for branching.
%In time $\mathcal{O}(nm)$, the sets $T_i^\nu,\, i \in \P$ can be computed.
%We distinguish two cases
%\todo{The exact choices here will need to be discussed}
%\begin{enumerate}[label=(\alph*)]
%    \item If there is no non-empty $T_i^\nu$, we choose an arbitrary $S \in \bs^\nu$ and split it fifty-fifty (or approximately)
%    \item Otherwise we choose a $T_i^\nu$ such that $S \cap T_i^\nu \neq \emptyset$ for some $S \in \bs^\nu$.
%    We choose $T_i^\nu$ such that $|S \cap T_i^\nu|$ is as large as possible.
%    Then .....
%\end{enumerate}

\subsection{Cutting planes}\label{subsec:ramp-cutting-planes}
As we have seen, the relaxation of \eqref{eq:RP} is not tight.
In practice, solving the MIP in this form is slow.
To tighten the formulation, we study cutting planes for \eqref{eq:RP-pos}.
Next, we will derive separation algorithms for both these cutting planes and the constraints \eqref{eq:RP-neg}.

When considering only the constraints \eqref{eq:RP-pos},
the problem is a special case of the Mixed Vertex Packing Problem~\cite{atamturk_mixed_2000}.
From this, it follows that for a given point $i$, the relevant cutting planes are
\begin{equation}\label{eq:RP-pos-cutting-planes}
\begin{split}
    &\gamma \leq 2 - \sum_{l\in T}\overline{\gamma}_{il}z_l, \quad \\
    &T = \{l_1, l_2, \dots l_t\} \subseteq \L \text{ such that } \gamma_{l_{j-1}} > \gamma_{l_{j}} \text { for } j = 2, \dots, t,
\end{split}
\end{equation}
where
\begin{align*}
    \overline\gamma_{il_1} &= 2 - \gamma_{il_1}, \\
    \overline\gamma_{il_j} &= 2 - \gamma_{il_j} - (2 - \gamma_{il_{j-1}}) = \gamma_{il_{j-1}} - \gamma_{il_j}, \quad j = 2, \dots, t.
\end{align*}
For more details we refer to \autoref{sec:cutting_planes_vertex_packing}.
Another way to approach the derivation of these cutting planes is via Extended Polymatroid Inequalities (EPIs),
recognizing that $f$ is supermodular.
\cite{kucukyavuz_mixed-integer_2023, atamturk_submodular_2020, lovasz_submodular_1983}
This approach is not treated in the report, but may be of interest to the reader.

The idea is now to solve \eqref{eq:RP} by adding the cutting planes \eqref{eq:RP-pos-cutting-planes} and the constraints
\eqref{eq:RP-neg} when necessary.
To do this, we need an efficient separation algorithm.


\subsection{Separation}\label{subsec:separation}
We now describe a polynomial-time separation algorithm for $\eqref{eq:RP-pos-cutting-planes}$ and then for $\eqref{eq:RP-neg}.$

Separation of \eqref{eq:RP-pos-cutting-planes} for fixed $i$ can be done in $\mathcal{O}(m  \log m)$ time using a greedy algorithm.
Without loss of generality, consider a fixed point $i$ and
\[
    2 \geq \gamma_1 \geq \gamma_2 \geq \cdots \geq \gamma_m \geq 0,
\]
each corresponding to a clause.
Assume the $z$-variables are ordered accordingly.
Let $(\bm{z}, \gamma)$ be a solution to the (partial) LP\@.
Note that $\bm{z} \geq 0$.
Observe that a violated constraint exists if and only if there exist
\[
    \tau_1 < \tau_2 < \cdots < \tau_t \in [m].
\] such that
\begin{equation}\label{eq:violated_pos_cons}
    \gamma > 2 - \sum_{\kappa = 1}^t \overline{\gamma}_{\tau_\kappa}z_{\tau_\kappa},
\end{equation}
where
\begin{align*}
    \overline\gamma_{\tau_1} &= 2 - \gamma_{\tau_1}, \\
    \overline\gamma_{\tau_j} &= 2 - \gamma_{\tau_j} - (2 - \gamma_{\tau_{j-1}}) = \gamma_{\tau_{j-1}} - \gamma_{\tau_j}, \quad j = 2, \dots, t.
\end{align*}
If the inequality holds for a minimizer of the right hand side, then we have found a new constraint.
If the inequality does not hold for this minimizer, we may conclude that no violated constraint exists.
Finding a minimizer of the right hand side is equivalent to finding a maximizer of
\begin{equation}\label{eq:sep-pos-sum}
    \sum_{\kappa = 1}^t \overline{\gamma}_{\tau_\kappa}z_{\tau_\kappa}.
\end{equation}

Denote $\gamma_0$ = 2 and $\tau_0 = 0$ for convenience.
We propose the following algorithm.

\begin{algorithm}[h!]
\caption{Algorithm to find a maximizer of \eqref{eq:sep-pos-sum}}
\label{alg:sep-pos}
\begin{algorithmic}[1]
\Require $\gamma, \bm{z}$
\State $i \gets 1$
\While{$\left\{\sigma \fw \gamma_\sigma < \gamma_{\tau_{i-1}}\right\} \neq \emptyset$}
    \State Choose a $\tau_i$ such that
    $z_{\tau_i} = \max\left\{z_\sigma : \gamma_{\sigma_i} < \gamma_{\tau_{i-1}}\right\}.$
    \State $i \gets i + 1.$
\EndWhile \\
\Return $(\tau_\kappa)_{\kappa=1,\dots,i-1}.$
\end{algorithmic}
\end{algorithm}

This algorithm can be implemented to run in $O(m\log m)$ time by first sorting (the indices of) $\bm{z}$
and iterating over the sorted list, keeping in mind that $\tau$ values must be increasing.

\begin{theorem}\label{thm:sep-pos}
    \autoref{alg:sep-pos} returns a subset \[\tau_1 < \tau_2 < \cdots < \tau_t \in [m]\] that maximizes
    \[
        \sum_{\kappa = 1}^t \overline{\gamma}_{\tau_\kappa}z_{\tau_\kappa}.
    \]
\end{theorem}

\begin{proof}
    Let
    \[
        \tau_1 < \tau_2 < \cdots < \tau_t \in [m]
    \]
    be a sequence.
    Define
    \[
        u_{\sigma} = \umax{\tau \geq \sigma} \; z_{\tau}, \tquad \sigma \in [m].
    \]
    Now observe that
    \begin{equation}\label{eq:sep-pos-proof-1}
        \begin{split}
            S(\tau) = \sum_{\kappa = 1}^t \overline{\gamma}_{\tau_\kappa}z_{\tau_\kappa}
            & = \sum_{\kappa = 1}^t z_{\tau_\kappa} \sum_{\sigma = \tau_{\kappa-1}+1}^{\tau_{\kappa}} (\gamma_{\sigma-1} - \gamma_{\sigma}) \\
            & = \sum_{\kappa = 1}^t \sum_{\sigma = \tau_{\kappa-1}+1}^{\tau_{\kappa}} z_{\tau_\kappa}(\gamma_{\sigma-1} - \gamma_{\sigma}) \\
            & \leq \sum_{\kappa = 1}^t \sum_{\sigma = \tau_{\kappa-1}+1}^{\tau_{\kappa}} u_{\sigma}(\gamma_{\sigma-1} - \gamma_{\sigma}) \\
            & \leq \sum_{\sigma = 1}^m u_{\sigma}(\gamma_{\sigma-1} - \gamma_\sigma).
        \end{split}
    \end{equation}
    The first inequality holds because
    \[
         \tau_\kappa \geq \sigma \quad \implies z_{\tau_\kappa} \leq u_\sigma,
    \]
    and $\gamma_{\sigma-1} - \gamma_\sigma \geq 0$ for all $\sigma \in [m]$.
    The second holds because the ranges $\{\tau_{\kappa-1}+1, \dots, \tau_\kappa\}, \; \kappa=1,\dots,t$
    partition $\{1, \dots, \tau_t\}$ and the omitted terms are non-negative.

    We now claim that the greedy method yields a solution that attains this upper-bound.
    Denote by $g_1 < g_2 < \dots < g_{t_g} \in [m]$ a solution produced by the greedy algorithm.
    Denote for convenience $g_0 = 0.$
    Let $\kappa \in [t_g].$ The candidate set in iteration $\kappa$ is
    \[
        \{\sigma \fw \gamma_{\sigma} < \gamma_{g_{\kappa-1}}\}.
    \]
    Since $\gamma$ values are non-increasing, it follows that this set equals
    \[
        \{\sigma \geq s_\kappa\}
    \]
    for some $s_\kappa > g_{\kappa-1}$.
    Hence, the choice in iteration $\kappa$, satisfies
    \[
        z_{g_\kappa} = u_{s_\kappa}.
    \]
    Now let $\sigma \in \{g_{\kappa-1} + 1, \dots, g_\kappa\}$ such that $\gamma_{\sigma} < \gamma_{\sigma-1} $.
    If that condition does not hold, then
    \[
        \gamma_{\sigma - 1} - \gamma_{\sigma} = 0,
    \]
    so $\sigma$ does not contribute to the bound derived in \eqref{eq:sep-pos-proof-1}.
    It follows that \[
                        \gamma_{\sigma} < \gamma_{\sigma-1} \leq \gamma_{g_{\kappa-1}}.
    \]
    So $\sigma \geq s_\kappa$.
    By definition of $u$, we then have that $u_{\sigma} \leq u_{s_\kappa} = z_{g_\kappa}.$
    Since $\sigma \leq g_\kappa$, by definition of $u$ it holds that $u_\sigma \geq z_{g_\kappa}.$
    So \[
           z_{g_\kappa} \leq u_{\sigma} \leq z_{g_\kappa}.
    \]
    Hence, equality holds on this entire range.

    For $\sigma > g_{t_g}, $ the greedy algorithm terminated, so there is no $\sigma'$ such that $\gamma_{\sigma'} < \gamma_{g_{t_g}}$,
    meaning $\gamma_\sigma = \gamma_{g_{t_g}}$ for all $\sigma > g_{t_g}$.
    Therefore
    \begin{equation}\label{eq:sep-pos-proof-2}
        \begin{split}
            S(\bm{g})
            & = \sum_{\kappa = 1}^{t_g} \sum_{\sigma = g_{\kappa-1}+1}^{g_\kappa}
                z_{g_\kappa}(\gamma_{\sigma-1} - \gamma_\sigma) \\
            & = \sum_{\kappa = 1}^{t_g} \sum_{\sigma = g_{\kappa-1}+1}^{g_\kappa}
                u_{\sigma}(\gamma_{\sigma-1} - \gamma_\sigma) \\
            & = \sum_{\sigma = 1}^m u_{\sigma}(\gamma_{\sigma-1} - \gamma_\sigma).
        \end{split}
    \end{equation}
    The second equality uses that $u_\sigma = z_{g_\kappa}$ whenever
    $\gamma_{\sigma-1} - \gamma_\sigma > 0$.
    The third uses that the ranges
    $\{g_{\kappa-1}+1, \dots, g_\kappa\}$ partition $\{1, \dots, g_{t_g}\}$
    and all remaining terms are zero.
    Combined with \eqref{eq:sep-pos-proof-1}, it follows that
    \[S(\bm{g}) \geq S(\bm{\tau})\]
    for any feasible sequence $\bm{\tau}$.
    So, $S(\bm{g})$ is a maximizer of \eqref{eq:sep-pos-sum}.
\end{proof}
\vspace{2em}
Next, consider $\bm{z} \in [0,1]^m$, the variables
\[
    \gamma_1 \leq \gamma_2 \leq \cdots \leq \gamma_m
\]
and the constraint family
\begin{equation}\label{eq:RP-neg-onedim}
    \gamma \geq \gamma_i - \sum_{i' = 1}^{i} (\gamma_i - \gamma_{i'}) z_{i'}, \tquad i \in [m],
\end{equation}
which translates directly to \eqref{eq:RP-neg}.
A violated inequality exists if and only if
\[
    \gamma < \gamma_i -  \sum_{i' = 1}^{i} (\gamma_i - \gamma_{i'}) z_{i'}
\]
for some $i \in [m].$
Therefore, we are interested in a maximizer of the right hand side.
We assume $m \geq 2$, as $m=1$ is a trivial case.
Define $g: [m] \to \R$ by
\[
    g(i) = \gamma_i - \sum_{i' = 1}^{i} (\gamma_i - \gamma_{i'}) z_{i'}.
\]
Let $i \in [m-1].$ Then
\[
    \begin{split}
        g(i+1) - g(i) &= \gamma_{i+1} - \sum_{i' = 1}^{i+1} (\gamma_{i+1} - \gamma_{i'}) z_{i'} - \gamma_i + \sum_{i' = 1}^{i} (\gamma_i - \gamma_{i'}) z_{i'} \\
        &= \gamma_{i+1} - \sum_{i' = 1}^{i} (\gamma_{i+1} - \gamma_{i'}) z_{i'} - \gamma_i + \sum_{i' = 1}^{i} (\gamma_i - \gamma_{i'}) z_{i'} \\ & \tquad -  (\gamma_{i+1} - \gamma_{i+1})z_{i+1}  \\
        &= \gamma_{i+1} - \gamma_i - \sum_{i' = 1}^{i} (\gamma_{i+1} - \gamma_{i'} - \gamma_i + \gamma_{i'}) z_{i'} \\
        &= \gamma_{i+1} - \gamma_i - \sum_{i' = 1}^{i} (\gamma_{i+1} - \gamma_i) z_{i'} \\
        &= (\gamma_{i+1} - \gamma_i) \left(1 - \sum_{i' = 1}^{i} z_{i'}\right).
    \end{split}
\]
Write
\[
    S_i = \sum_{i'=1}^{i} z_{i'},
\] so that
\[
    g(i+1) - g(i) = (\gamma_{i+1} - \gamma_i)(1 - S_i).
\]
Because $\bm{z} \geq 0$, the sequence $(S_i)_{i \in [m]}$ is non-decreasing.
Since $\gamma_{i+1} \geq \gamma_i$ for all $i \in [m-1]$, it follows that
\begin{equation}\label{eq:sep-neg-obs-1}
    g(i+1) > g(i) \quad \iff \quad  S_i < 1 \quad \wedge \quad  \gamma_{i+1} > \gamma_i.
\end{equation}
Furthermore,
\begin{equation}\label{eq:sep-neg-obs-2}
    S_i \leq 1 \quad \implies g(i+1) \geq g(i).
\end{equation}
\begin{theorem}\label{thm:sep-neg-opt}
    Let $i^* \in [m-1]$ be the largest index such that
    \[
        S_{i^*} = \sum_{i' = 1}^{i^*} z_{i'} \leq 1,
    \]
    which exists because $S_1 = z_1 \leq 1$.
    Then $i^*+1$ is a maximizer for $g$.
\end{theorem}

\begin{proof}
    Let $i^*$ be as described.
    Recall that $(S_i)_{i \in [m]}$ is non-decreasing.
    We have that $S_i \leq 1$ for $i \leq i^*$ and $S_i > 1$ for $i \in [m]$ such that $i > i^*$.
    It follows by \eqref{eq:sep-neg-obs-2} that $g(i+1) \geq g(i)$ for all $i \in [m-1]$ such that $i \leq i^{*}.$
    By the contrapositive of the forward implication in \eqref{eq:sep-neg-obs-1}, we furthermore have that
    $g(i+1) \leq g(i)$ for every $i \in [m-1]$ such that $i \geq i^*+1$.
    Therefore, $g(i) \leq g(i^*+1)$ for all $i \in [m]$.
    We conclude that $i^*+1$ is a maximizer for $g$.
\end{proof}
Consider $i^*$ as described by \autoref{thm:sep-neg-opt}.
If $g(i^*+1) > \gamma$, then the corresponding constraint \eqref{eq:RP-neg-onedim} is violated.
Otherwise, no violated inequality exists.
Finding $i^*$ can be done in linear time.
Sorting $\bm{\gamma}$ is done during pre-processing.

\section{Examples}\label{sec:ramp-examples}
In this section, we study the behavior of the ramp loss model on some simple, artificial data.
These examples employ a ``full master problem'', meaning all relevant rules (columns) are considered.
Precisely, given a positive integer $D$, all possible rules with at most $D$ clauses satisfying the following conditions are considered:
\begin{itemize}
    \item A feature appears at most once in any rule.
    So a feature appears in at most one clause within a rule.
    \item A feature and their negated version do not appear together in any rule.
    This means a feature cannot appear in both a ``$\geq$''-clause and a ``$\leq$''-clause within a rule.
\end{itemize}
The examples are in two dimensions.
Recall that for $\beta = 0$, the ramp loss model is equivalent to the threshold model.

\begin{example}[Two clouds]\label{ex:two_clouds}
    \begin{figure}
        \includesvg[width=\textwidth]{gen/test_clouds/example_1.svg}
        \caption{Instance of \autoref{ex:two_clouds}. The example consists of a negative and a positive cloud of $100$ points each, following $\mathcal{N}((8, 10), 2)$ and $\mathcal{N}((12, 10), 2)$ respectively.}
        \label{fig:two_clouds}
    \end{figure}
    We construct two clouds of points according to multivariate normals.
    Precisely,
    \begin{itemize}
        \item A cloud of $100$ negative points drawn from $\text{Normal}((8, 10), 2)$,
        \item A cloud of $100$ positive points drawn from $\text{Normal}((12, 10), 2)$.
    \end{itemize}
    Hence, the cloud centers share $x_2$ but differ in $x_1$.
    An instance is depicted in \autoref{fig:two_clouds}
    Choosing $c=1$ as rule complexity, we solve the ramp loss model for the parameters
    \[C \in \{1,2\}, \quad \beta \in \{0, 0.5, 1, 2, 3, 4, 5\}, \quad D = \{2\}.\]

    \begin{figure}
        \includesvg[width=.8\textwidth]{gen/test_clouds/results/O_1_C_1_mw_0_additive.set.svg}
        \caption{Solution to an instance of \autoref{ex:two_clouds} computed with the ramp loss model, configured with $C=1$, $\beta=0$, $D=2$.}
        \label{fig:two_clouds_solution_simple}
    \end{figure}
    For $\beta=0,$ the solution contains only the rule
    $x_1 \geq 9.88$ regardless of $C$.
    So, for this example, an increase in $C$ does not cause the model to behave differently.
    This solution is depicted in \autoref{fig:two_clouds_solution_simple}.

    For $C=1$ fixed, the same solution is found for $\beta$ up until and including $\beta = 4$.
    For $C=1, \; \beta=5,$ the rule is $x_1 \geq 10.73$ instead.
    We conclude that for $C=1$, the ramp loss model finds a sensible solution for any of the tested $\beta$.

    For $C > 1$, $\beta > 0$, the model behaves peculiarly: all solutions consist of two rules.
    For $C=2$, $\beta = 0.5$ the ruleset consists of two adjacent, non-overlapping rules, as depicted in \autoref{fig:two_clouds_solution_bad_a}.
    The same holds for $C=2$, $\beta=1$, as shown in \autoref{fig:two_clouds_solution_bad_b}.
    These rulesets achieve a much lower objective than the rule $x_1 \geq 9.88$ found earlier.
    For $\beta=0.5,$ this can be partly attributed to the exclusion of the lower left corner in the covered region.
    More interestingly, for $\beta=1$, the choice of this ruleset can be attributed to the additive nature of the model.
    To see this, recall that points near a margin induce a loss, even if the point resides on the desired side of the covered region.
    Looking closer, constraint \eqref{eq:RM-misP} implies that the loss of a positive point decreases as the number of margins it is in increases.
    In particular, being in two or more margins can be equivalent to being far inside the covered region, which gives a loss of $0$.
    This encourages the use of more than one rule.
    From $\beta=2$ onwards, the ruleset consists of two rules, one of which is contained in the other.
    This is depicted in \autoref{fig:two_clouds_solution_bad_2}.
    The objective is again much lower than the simple rule $x_1 \geq 9.88$.
    This phenomenon can again be attributed to the additive nature of the model.

    \begin{figure}
        \centering
        \begin{subfigure}{.48\textwidth}
            \centering
            \includesvg[width=\textwidth]{gen/test_clouds/results/O_1_C_2_mw_0.5_additive.set.svg}
            \caption{}
            \label{fig:two_clouds_solution_bad_a}
        \end{subfigure}
        \hfill
        \begin{subfigure}{.48\textwidth}
            \centering
            \includesvg[width=\textwidth]{gen/test_clouds/results/O_1_C_2_mw_1_additive.set.svg}
            \caption{}
            \label{fig:two_clouds_solution_bad_b}
        \end{subfigure}
        \caption{Solutions to an instance of \autoref{ex:two_clouds} computed with the ramp loss model, configured with $C=2$, $\beta=0.5$ (a) and $C=2$, $\beta=1$ (b), respectively.}
        \label{fig:two_clouds_solution_bad}
    \end{figure}

    \begin{figure}
        \centering
        \begin{subfigure}{.48\textwidth}
            \centering
            \includesvg[width=\textwidth]{gen/test_clouds/results/O_1_C_2_mw_2_additive.set.svg}
            \caption{}
            \label{fig:two_clouds_solution_bad_2_a}
        \end{subfigure}
        \hfill
        \begin{subfigure}{.48\textwidth}
            \centering
            \includesvg[width=\textwidth]{gen/test_clouds/results/O_1_C_2_mw_3_additive.set.svg}
            \caption{}
            \label{fig:two_clouds_solution_bad_2_b}
        \end{subfigure}
        \caption{Solutions to an instance of \autoref{ex:two_clouds} computed with the ramp loss model, configured with $C=2$, $\beta=2$ (a) and $C=2$, $\beta=3$ (b), respectively.}
        \label{fig:two_clouds_solution_bad_2}
    \end{figure}

\end{example}

\begin{example}[Three clouds]\label{ex:three_clouds}
    \begin{figure}
        \includesvg[width=\textwidth]{gen/test_clouds/example_2.svg}
        \caption{Instance of \autoref{ex:three_clouds}. The example consists of a negative cloud of $100$ points, a positive cloud of $100$ points and a positive cloud of $10$ points.
        The point clouds have distributions $\text{Normal}((8, 10), 2)$, $\text{Normal}((12, 10), 2)$ and $\text{Normal}((5,15), 1)$, respectively.}
        \label{fig:three_clouds}
    \end{figure}
    Consider \autoref{ex:two_clouds} with an additional positive cloud of $10$ points, drawn from $\text{Normal}((5, 15), 1)$.
    An instance is depicted in \autoref{fig:three_clouds}.
    We again solve the ramp loss model for parameters
    \[C \in \{1,2,3\}, \quad \beta \in \{0, 0.5, 1, 2, 3, 4, 5\}, \quad D\in \{2\}.\]
    For $C = 1$, the rule $x_1 \geq 9.71$ is found for all $\beta$ up until and including $\beta = 4$.
    For $C = 1$, $\beta = 5$, the rule is $x_1 \geq 8.77$ instead.
    Notably, the margin covers every point in the positive cloud of $10$ points in this case.
    For $C = 2$ and $\beta \in \{0, 0.5, 1\},$ the ruleset
    \[
        x_1 \geq 9.71 \quad \vee \quad  \left(x_1 \leq 8.03 \; \wedge \; x_2 \geq 13.36\right)
    \]
    is found.
    This solution is depicted in \autoref{fig:three_clouds_solution_good}.
    \begin{figure}
        \includesvg[width=.8\textwidth]{gen/test_clouds/results/O_2_C_2_mw_0_additive.set.svg}
        \caption{Solution to an instance of \autoref{ex:three_clouds} computed with the ramp loss model, configured with $C=2$, $\beta=0$, $D=2$.}
        \label{fig:three_clouds_solution_good}
    \end{figure}
    For $C=2$, $\beta=2$ the resulting ruleset is different and may be considered reasonable.
    This ruleset is depicted in \autoref{fig:three_clouds_solution_bad_a}.
    For $C=2$, $\beta>2$, rulesets consist of overlapping rules and ignore the positive cloud of 10 points.
    An example is depicted in \autoref{fig:three_clouds_solution_bad_b}.
    For $C=3$, the solutions are similar to the $C=2$ case of \autoref{ex:two_clouds},
    but with an extra rule to accommodate the additional positive point cloud.
    An example is shown in \autoref{fig:three_clouds_solution_bad_2}.

    \begin{figure}
        \centering
        \begin{subfigure}{.48\textwidth}
            \centering
            \includesvg[width=\textwidth]{gen/test_clouds/results/O_2_C_2_mw_2_additive.set.svg}
            \caption{}
            \label{fig:three_clouds_solution_bad_a}
        \end{subfigure}
        \hfill
        \begin{subfigure}{.48\textwidth}
            \centering
            \includesvg[width=\textwidth]{gen/test_clouds/results/O_2_C_2_mw_3_additive.set.svg}
            \caption{}
            \label{fig:three_clouds_solution_bad_b}
        \end{subfigure}
        \caption{Solutions to an instance of \autoref{ex:three_clouds} computed with the ramp loss model, configured with $C=2$, $\beta=2$ (a) and $C=2$, $\beta=3$ (b), respectively.}
        \label{fig:three_clouds_solution_bad}
    \end{figure}

    \begin{figure}
        \includesvg[width=.8\textwidth]{gen/test_clouds/results/O_2_C_3_mw_1_additive.set.svg}
        \caption{Solution to an instance of \autoref{ex:three_clouds} computed with the ramp loss model, configured with $C=3$, $\beta=1$, $D=2$.}
        \label{fig:three_clouds_solution_bad_2}
    \end{figure}

\end{example}

In these two examples, the model behaves as desired when $\beta = 0$.
When a margin is introduced, freedom to choose rules results in undesirable rulesets.
The clearest example of this is \autoref{fig:three_clouds_solution_bad_b}, where a positive point cloud is ignored,
because overlapping rules are beneficial in terms of loss, due to additivity in the model.
From a practitioner's perspective, the problem is sometimes solved by limiting $C$ and $\beta$.
After all, for adequate choice of $C$ and $\beta$, the model produces results considered reasonable.
However, the results suggest that the model in its current form is not suitable for the problem at hand.
Therefore, we introduce a version of the model where positive points no longer benefit from being in more than one margin.
We refer to this model as the non-additive ramp loss moel.

\section{Non-additive variant}\label{sec:ramp-nonadditive}
In the non-additive ramp loss model, we wish to replace
\[
    \zeta_i + \sum_{k \in \K} \gamma_{ik} w_k \geq 2, \tquad i \in \P,
\]
by
\[
    \zeta_i + \umax{k \in \K} \; \gamma_{ik} w_k \geq 2, \tquad i \in \P.
\]
which is equivalent to
\begin{equation*}
    \begin{split}
        \zeta_i - 2 &\geq - \umax{k \in \K} \; \gamma_{ik} w_k\\
        &= \umin{k \in \K} \; -\gamma_{ik} w_k , \tquad i \in \P. \\
    \end{split}
\end{equation*}
Define
\[
    \eta_{ik} = -\gamma_{ik}, \tquad i \in \P, \; k \in \K.
\]
Then one formulation of the non-additive ramp-loss model is
\begin{subequations}
    \label{eq:NRM}
    \begin{align}
        \mathbf{\min} \quad & \sum_{i \in \cP} \zeta_i + \sum_{i \in \cN} \sum_{k \in \K} \gamma_{ik} w_k \label{eq:NRM-obj} \\
        \mathbf{\st} \quad & \zeta_i - 2 \geq \eta_{ik'} - \sum_{\substack{k \in \K, \\ \eta_{ik} < \eta_{ik'}}} (\eta_{ik'}-\eta_{ik}) w_k,
        & \forall i \in \cP, \; k' \in \K, \label{eq:NRM-misP} \\
        & \zeta_i - 2 \geq - \sum_{k \in \K} 2w_k, & \forall i \in \I, \label{eq:NRM-fix}\\
        & \sum_{k \in \K} c_k w_k \leq C, & \label{eq:NRM-complex} \\
        & w \in \{0, 1\}^{\K}, \zeta \in [0,2]^{\cP}. \label{eq:NRM-vars}
    \end{align}
\end{subequations}
The reduced cost of a rule is derived in \autoref{sec:reduced_costs_nonadditive_ramp_loss}.
However, this reduced cost cannot be easily be modeled inside a MIP\@.

Alternatively, we may consider the following formulation:
\begin{subequations}
    \label{eq:NRMa}
    \begin{align}
        \mathbf{\min} \quad & \sum_{i \in \cP} \zeta_i + \sum_{i \in \cN} \sum_{k \in \K} \gamma_{ik} w_k \label{eq:NRMa-obj} \\
        \mathbf{\st} \quad & \zeta_i - 2 \geq \sum_{k \in \K} u_{ik}\eta_{ik}, & \forall i \in \cP, \label{eq:NRMa-misP} \\
        & \sum_{k \in \K} u_{ik} = 1, & \forall i \in \cP, \label{eq:NRMa-misP2} \\
        & u_{ik} \leq w_k, & \forall i \in \cP, \; k \in \K, \label{eq:NRMa-misP3} \\
        & \sum_{k \in \K} c_k w_k \leq C, & \label{eq:NRMa-complex} \\
        & w \in \{0, 1\}^{\K}, \zeta \in [0,2]^{\cP}, u \in [0,1]^{\P \times \K}. \label{eq:NRMa-vars}
    \end{align}
\end{subequations}
While this formulation has a large number of variables, the $u$-variables need not be binary.
The reduced cost of a variable $w_k,\; k \in \K$, is derived in \autoref{sec:reduced_costs_nonadditive_ramp_loss}.
Implementing column-and-row generation for this master problem formulation is not trivial,
as a rule corresponds to a $w$-variable and many $u$-variables.

\begin{note}
    A formulation similar to \eqref{eq:NRMa} may be constructed for the pricing problem \eqref{eq:RP_set_formulation}.
    Implementing and solving this MIP for a number of test cases revealed that it is not easily solved by SCIP.\@
    Using proprietary solver Gurobi, as well as some tricks exploiting the nature of the data,
    performance improved moderately, but not enough to solve the problem in a handful of minutes.
    Hence, the more compact formulation \eqref{eq:RP} was explored in combination with separation of cutting planes.
    Ultimately, none of the approaches resulted in the ability to solve the pricing problem efficiently for our dataset.
\end{note}

Evidently, solving the relaxed non-additive ramp loss model exactly is difficult,
as even the LP cannot be easily solved via column (and row) generation.
A pragmatic approach is to rely solely on a pricing heuristic.
Adapting the heuristic pricer to the non-additive model is possible, but requires substantial work.
Therefore, we first explore whether the non-additive model gives better results on example instances.
We consider again the setting and examples in \autoref{sec:ramp-examples}.
We use formulation \eqref{eq:NRM}, as it is easiest to implement within the existing implementation of the ramp loss model.

\begin{note}
    In the construction and implementation of \eqref{eq:NRM}, constraint \eqref{eq:NRM-fix} was overlooked.
    This affects the objective of the empty ruleset if there exist $i \in \P$ such that $\eta_{ik} < 0$ for all
    $k$ considered in the branching phase.
    The objective of a non-empty ruleset is correctly modeled.
    Since we are interested in an empty ruleset, and non-empty rulesets are indeed found as solutions,
    the mistake of omitting \eqref{eq:NRM-fix} seems to have no effect in practice.
\end{note}

\begin{example}[Two clouds, non-additive model]\label{ex:two_clouds_nonadditive}
    Consider again \autoref{ex:two_clouds} and the tested parameters.

    For $C=1$ and any $\beta$, the results are identical to those of the additive model.
    For $C=2,$ $\beta=0,$ the ruleset differs and is depicted in \autoref{fig:two_clouds_nonadditive_solution_bad_a}.
    It consists of two rules, and the covered region is identical to the one-rule solution found for the additive model.
    The objective is also identical.
    To the model, both solutions are equally good.
    For classification tasks, they are also identical.
    If the smaller ruleset is preferred, for instance because of interpretability, then $C$ could be chosen as low as possible so long as the performance,
    whichever way it is measured, does not decrease.
    Another solution could be to penalize ruleset complexity in the model.
    It is important to note that this solution, when considered in the additive model with $\beta =0$,
    has the exact same objective value as \autoref{fig:two_clouds_solution_simple}.
    Hence, the choice of this solution in the non-additive model is a coincidence rather than an artifact.
    In both the additive and non-additive model, the value of $C$ is best kept small to prevent this.

    For $C=2,$ $\beta=0.5$ the solution is similar to \autoref{fig:two_clouds_solution_bad_a} and is shown in \autoref{fig:two_clouds_nonadditive_solution_bad_b}.
    For $C=2,$ $\beta=1$, the solution contains $x_1 \geq 9.88$, but also another rule, as shown in \autoref{fig:two_clouds_nonadditive_solution_bad_c}.
    Removing this other rule from the ruleset, leaving just the rule $x_1 \geq 9.88$, results in the same objective.
    Hence, one may want to prune solutions.
    One way of pruning solutions is to consider all subsets of the ruleset, and pick the smallest subset that has the same objective value as the full ruleset.
    Alternatively, one could again limit $C$ or penalize rule complexity.

    Notably, for $\beta > 0.5$, the non-additive model seems less included to produce undesirable rulesets than the additive model.
    While the additive model yields two-rule solutions for $C=2,$ $\beta>0.5,$
    the non-additive yields a one-rule solution either immediately or after pruning.
    This can be partly attributed to the non-additivity that was introduced.
    Another factor is the negative points' loss, which is calculated in the same way in either model.
    The loss for negative points is additive, hence using more rules is bad for negative points in or around the covered region.

    \begin{figure}
    \centering
    \begin{subfigure}{.48\textwidth}
        \centering
        \includesvg[width=\textwidth]{gen/test_clouds/results/O_1_C_2_mw_0_nonadditive.set.svg}
        \caption{}
        \label{fig:two_clouds_nonadditive_solution_bad_a}
    \end{subfigure}
    \hfill
    \begin{subfigure}{.48\textwidth}
        \centering
        \includesvg[width=\textwidth]{gen/test_clouds/results/O_1_C_2_mw_0.5_nonadditive.set.svg}
        \caption{}
        \label{fig:two_clouds_nonadditive_solution_bad_b}
    \end{subfigure} \\
        \begin{subfigure}{.48\textwidth}
        \centering
        \includesvg[width=\textwidth]{gen/test_clouds/results/O_1_C_2_mw_1_nonadditive.set.svg}
        \caption{}
        \label{fig:two_clouds_nonadditive_solution_bad_c}
    \end{subfigure} \\
    \caption{Solutions to an instance of \autoref{ex:two_clouds} computed with the ramp loss model, configured with $C=2$, $\beta=0.5$ (a) and $C=2$, $\beta=1$ (b), respectively.}
    \label{fig:two_clouds_nonadditive_solution_bad}
\end{figure}

\end{example}

We conclude that the non-additive ramp loss model behaves better than the additive ramp loss model.
In practice, we perform pricing on the additive variant.
Then, we use the non-additive variant inside the branching stage.
This effectively means we take the columns generated using the additive variant,
add these columns to an IP of the non-additive variant, and solve.

\section{Implementation}\label{sec:ramp-implementation}
The custom solver described in \autoref{sec:boolean-implementation} is capable of solving the ramp loss model.
This means there is another option, namely the model type, which is either ``Boolean'' or ``Threshold''.
The model type ``Threshold'' solves both the threshold model (\autoref{ch:threshold-model})
and the ramp loss model, as the threshold model is a special case of the ramp loss model,
precisely the case $\beta = 0.$

The implementation for both model types share relevant parts of the codebase.
Notably, both model types use the same basis for the heuristic pricer.
The key difference is how the reduced cost is computed inside the search tree.
We will now discuss how this is done for the ramp loss model.
Let $k_\text{parent}, k_\text{child}$ be the rules corresponding to some parent and child node, respectively.
Then the reduced cost of the child node is given by
\begin{equation}
        \hat \rho_{k_\text{child}} =  \sum_{i \in \cN} \gamma_{ik_\text{child}}
    -\sum_{i \in \cP}  \mu_i \gamma_{ik_\text{child}} +\lambda  c_{k_\text{child}}.  \label{eq:RP_redcost_node}
\end{equation}
We now wish to express $\gamma_{ik_\text{child}}$ in terms of $\gamma_{ik_\text{parent}}$.
Ideally, we can then store $\gamma$-values cleverly in the search tree, and compute
$\hat \rho_{k_\text{child}}$ from $\hat \rho_{k_\text{parent}}$ with relatively little effort.
Denote by $l$ the clause that is added in $k_\text{child}$ with respect to $k_\text{parent}.$
Recall the definition of $\gamma_{ik}$ in \eqref{eq:f_min}.
Evidently,
\[
  \gamma_{ik_\text{child}} \leq \gamma_{ik_\text{parent}}.
\]
Precisely, observe that
\begin{equation}\label{eq:RP_gamma_diff}
    \gamma_{ik_\text{child}} = \gamma_{ik_\text{parent}} - (\gamma_{ik_\text{parent}} - \gamma_{il})\Ind\{\gamma_{il} < \gamma_{ik_\text{parent}}\}.
\end{equation}
For convenience, we denote
\begin{equation}\label{eq:RP_delta_gamma}
    \Delta \gamma_{ik} = (\gamma_{ik_\text{parent}} - \gamma_{il})\Ind\{\gamma_{il} < \gamma_{ik_\text{parent}}\}
\end{equation}
It follows that
\begin{equation}
    \begin{split}
        \hat \rho_{k_\text{child}} = \; & \hat{\rho}_{k_\text{parent}} - \sum_{i \in \cN} \Delta \gamma_{ik_\text{child}}
        + \sum_{i \in \P} \Delta \mu_i \gamma_{ik_\text{child}}
        + \lambda.
    \end{split}
    \label{eq:RP_redcost_node_deriv}
\end{equation}
There are two important cases in which $\Delta \gamma_{ik} = 0$.
First, if $\gamma_{ik_\text{parent}} = 0,$ as $\gamma_{il}$ cannot be lower.
Second, if $\gamma_l = 2,$ as $\gamma_{ik_\text{parent}}$ cannot be higher.
Define
\begin{equation}
    \overline \I_k \coloneqq \{i \in \I \fw \gamma_{ik} > 0\} \tquad k \in \K.\label{eq:ramp_I_k}
\end{equation}
In other words, $\overline \I_k$ is the set of points that are covered by rule $k$ or are in the margin.
Define $\overline \N$ and $\overline \P$ analogously.
Then, instead of summing over $\N$ and $\P$, we may sum over
\[
    \overline \N \cap \{i \in \N : \gamma_{il} < 2\} \quad \text{and} \quad \overline \P \cap \{i \in \P : \gamma_{il} < 2\},
\]
respectively.
Observe that
\[
    \overline \I_{k_\text{child}} \subseteq \overline \I_{k_\text{parent}}.
\]
and
\[
    \overline \I_{k_\text{child}} = \overline \I_{k_\text{parent}} \setminus \{i \in \I \fw \gamma_{il} = 0\}.
\]
We store $\overline \I_k$ for each node $k$ in the search tree.
Furthermore, we do not store every $\gamma_{ik}$, as this would require many variables in each node of the search tree.
Instead, we store the key-value pair $(i, \gamma_{ik_\text{child}})$ at node $k_\text{child}$ if $\gamma_{ik_\text{child}}$ changed with respect to the parent.
Hence, finding the value of $\gamma_{ik_\text{parent}}$ requires traversing up the search tree until there is a hit.

Finally, we note that allowing a TF to appear at most once in a rule can speed up heuristic pricing.
When this restriction is enforced, the number of possible children of each node is reduced.
This is relatively easy to implement for the threshold model and the ramp loss model,
but is less convenient to implement for the boolean model, as it would require additional input to specify which binarized features
belong to the same TF\@.

\section{Methodology}\label{sec:ramp-methodology}
We conduct the experiments described in \autoref{sec:threshold-methodology} again,
but now the non-additive variant for various margin widths $\beta > 0.$
In particular, we consider \[
                               \beta=0, 0.2, 0.4, 0.6, 0.8, 1.
\]
We choose the range $0$ to $1$ because the standard deviation of most TF is around $2$.
We compare these results with the earlier results, which will be indicated by $\beta=0.$

We do not expect train accuracy to improve.
After all, the model directly maximizing accuracy is the original master problem \eqref{eq:BM_orig}.
Now we are considering a new variant, based on the Hamming loss, in an attempt to obtain qualitatively better results,
but not necessarily better train accuracy.

Instead, we will look at the qualitative results, in particular the ruleset frequency under cross-validation.
We also look at generalization (train and test accuracy).
We aim for a better generalization, i.e., a higher test accuracy.



\section{Results}\label{sec:ramp-results}
The train and test accuracies are given in \autoref{fig:pdl1_experiment_2_cross_train_test_mw}.
It shows results for $Q_y=0.5$ (top row) and $Q_y=0.6$ (bottom row).
The first column corresponds to $\beta=0$.
Hence, the first column is a variant of \autoref{fig:threshold_train_test_accuracy_mw_0.svg}.
The margin width $\beta$ increases with the column index, the last column corresponding to $\beta=1$.

From this plot, there is no clear improvement in test accuracy for $\beta>0$.
Increasing $\beta$ does appear to dampen the effect of $C$ on train accuracy.
For larger $\beta$, increasing $C$ leads to less improvement in train accuracy.
Hence, the ramp loss model may be less prone to overfitting.
It may also indicate that the ramp loss model is less capable of learning altogether.
Considering the confidence intervals for test accuracy,
no significant improvement in test accuracy can be established.
There may even be a decrease in test accuracy as $\beta$ increases.

Based on the plot, we study $\beta=0.6$ in more detail.
The ruleset frequency table is given in \autoref{tab:threshold_cross_interpretable_mw_0.6_C_1}.
For $Q_y=0.5$, the ruleset found is now consistent among all splits.
The opposite is true for $Q_y=0.6$.
Interestingly, the reverse was true for $\beta=0$, as seen in \autoref{tab:threshold_cross_interpretable_mw_0_C_1}.
It appears therefore that the introduction of the ramp loss,
which lead to the development of the non-additive model, did not lead to more consistent rulesets being found.
Hence, there appears to be no advantage of using the ramp loss model over the threshold model.

\clearpage

\begin{landscape}
    \vspace*{\fill}
    \begin{figure}[h]
        \includesvg[width=\textwidth]{gen/PDL1ex2cross/train_test_mw.svg}
        \caption{Confidence intervals for train and test accuracies for the ramp loss model, $C=1,2,3$, $D=2$, \\ $\beta=0,0.2,0.4,0.6,0.8,1$.}
        \label{fig:pdl1_experiment_2_cross_train_test_mw}
    \end{figure}
    \vspace*{\fill}
\end{landscape}


\begin{table}[h]
    \centering
    \begin{tabular}{llllr}
\toprule
 &  &  &  & \textbf{Frequency} \\
$\mathbf{Q_y}$ & $\mathbf{C}$ & $\mathbf{D}$ & \textbf{TF combination} &  \\
\midrule
\multirow[t]{2}{*}{0.5} & \multirow[t]{2}{*}{1} & 2 & {RELA $\wedge$ STAT2} & 10 \\
\cline{3-5}
 &  & 3 & {FOS $\wedge$ RELA $\wedge$ STAT2} & 10 \\
\cline{1-5} \cline{2-5} \cline{3-5}
\multirow[t]{6}{*}{0.6} & \multirow[t]{6}{*}{1} & \multirow[t]{2}{*}{2} & {RELA $\wedge$ STAT2} & 5 \\
 &  &  & {FOSL1 $\wedge$ RELA} & 5 \\
\cline{3-5}
 &  & \multirow[t]{4}{*}{3} & {FOS $\wedge$ RELA $\wedge$ STAT2} & 4 \\
 &  &  & {FOSL1 $\wedge$ RELA $\wedge$ STAT2} & 3 \\
 &  &  & {FOSL1 $\wedge$ IRF1 $\wedge$ RELA} & 2 \\
 &  &  & {FOS $\wedge$ FOSL1 $\wedge$ RELA} & 1 \\
\cline{1-5} \cline{2-5} \cline{3-5}
\bottomrule
\end{tabular}

    \caption{}
    \label{tab:threshold_cross_interpretable_mw_0.6_C_1}
\end{table}


